\section*{Capítulo \ref{ch:g1:living-or-not} --- Vivo ou não vivo?}

\begin{solution}{exo:g1:living-or-not:1}
Nasce, se alimenta, cresce, pode ter filhotes, morre. Hoje dá para
ver um gatinho se alimentando (mamando, comendo) e crescendo --- ele
fica maior semana após semana.
\end{solution}

\begin{solution}{exo:g1:living-or-not:2}
(a) Viva --- ela saiu de um ovo, se alimenta de néctar e pode ter
lagartinhas filhotes. (b) Não viva --- montada numa fábrica, nunca
se alimenta nem cresce. (c) Viva --- ela cresceu de uma semente.
(d) Não viva --- a nuvem cresce e diminui, mas nunca nasceu, não se
alimenta e não tem filhotes.
\end{solution}

\begin{solution}{exo:g1:living-or-not:3}
Ela não está viva agora. Já esteve viva: cresceu na árvore, e a
árvore a alimentava. Ela vem de um \omterm{def:g1:living-or-not:living}{ser vivo} --- a árvore.
\end{solution}

\begin{solution}{exo:g1:living-or-not:4}
Faça as perguntas da vida todas juntas: o rio nasceu? Ele se
alimenta? Ele cresce? Ele pode ter riozinhos filhotes? Não --- se
mexer e fazer barulho não são sinais de vida sozinhos. O rio não é
vivo, mesmo que muitos \omterm{def:g1:living-or-not:living}{seres vivos} morem \emph{dentro} dele.
\end{solution}

\begin{solution}{exo:g1:living-or-not:5}
Vivos: a aranha, a maçã na árvore. Já foram vivos (ou vêm de um \omterm{def:g1:living-or-not:living}{ser
vivo}): a meia de lã (lã de ovelha), a pinha caída (ela cresceu num
pinheiro). Nunca foi viva: a gota de chuva.
\end{solution}

\begin{solution}{exo:g1:living-or-not:6}
Não. A chama não passa nas perguntas da vida: ela foi acesa, não
nasceu, e, acima de tudo, não existem chamas filhotes que crescem
--- uma chama pode ser copiada para outra vela, mas não tem
filhotes como um caracol ou uma margarida. (É uma boa pergunta: a
chama engana mesmo dois dos sinais!)
\end{solution}

\begin{solution}{exo:g1:living-or-not:7}
Por exemplo: uma cadeira de madeira (de uma árvore), uma camiseta de
algodão (de um algodoeiro), um cobertor de lã (de uma ovelha), um
sapato de couro (de uma vaca), o mel da prateleira (das abelhas).
\end{solution}

\begin{solution}{exo:g1:living-or-not:8}
O feijão muda: ele incha, depois sai uma raizinha branca, depois um
broto verde. A pedrinha continua exatamente igual. Isso mostra que o
feijão estava vivo --- um \omterm{def:g1:living-or-not:living}{ser vivo} quietinho, esperando --- e que a
pedrinha não está.
\end{solution}

\begin{solution}{exo:g1:living-or-not:9}
Nasceu? O caracol saiu de um ovo; o robô foi montado. Pode ter
filhotes? Existem caracóis filhotes; robôs filhotes não existem.
(A alimentação também serve: a alface faz o caracol crescer, a pilha
não faz o robô crescer.)
\end{solution}

\begin{solution}{exo:g1:living-or-not:10}
O cogumelo nasce (ele aparece de coisinhas minúsculas na terra), se
alimenta (tira o alimento da terra e das folhas caídas), cresce ---
às vezes numa única noite ---, faz cogumelos novos e morre. Todos os
sinais da vida estão lá: ele é um \omterm{def:g1:living-or-not:living}{ser vivo}, mesmo sem pernas e sem
folhas verdes.
\end{solution}
